\chapter{Anexos de trazabilidad y reproducibilidad}
\begingroup
\sloppy

\section{Matriz objetivo--evidencia}

La Tabla~\ref{tab:matriz-objetivos-anexo} resume el estado verificable con los cuatro objetivos suministrados el 7 de
septiembre de 2026, actualizado con la fase física hasta el 21 de octubre de 2026. ``Parcial'' no significa incumplimiento oculto: identifica
el producto disponible y la evidencia que todavía falta para cerrar el
objetivo en el prototipo físico.

\begin{longtable}{p{0.09\textwidth}p{0.27\textwidth}p{0.31\textwidth}p{0.22\textwidth}}
\caption{Matriz resumida de objetivos, evidencia y estado.}\label{tab:matriz-objetivos-anexo}\\
\toprule
Objetivo & Resultado demostrable & Evidencia principal & Estado y faltante \\
\midrule
\endfirsthead
\toprule
Objetivo & Resultado demostrable & Evidencia principal & Estado y faltante \\
\midrule
\endhead
OE1 & Doce articulaciones ejecutadas en el robot suspendido, con topes, límite de cambio y realimentación de diez. & Capítulo~\ref{cap:plataforma-fisica}, registros de caminatas, firmware y nueve escenarios del supervisor. & Parcial; seguridad eléctrica, alimentación separada y migración del bucle a la Raspberry~Pi. \\
OE2 & Modelo nominal con geometría medida y actuador identificado. & Capítulo~\ref{cap:modelo}, \path{nova_sm3_real.xml} y modelos del actuador. & Parcial; identificar masas, inercias y fricción. \\
OE3 & Políticas PPO residuales entrenadas con el actuador real y seleccionadas en el robot; mejora del 15--19~\% en gateo. & Capítulo~\ref{cap:ppo-fisico}, Tabla~\ref{tab:ciclo-mejora} y políticas en \path{ppo_v10_robot/}. & Parcial; efecto sobre la locomoción apoyada sin medir. \\
OE4 & Seguimiento articular y repetibilidad medidos en el robot suspendido; margen nominal en simulación. & Diagnósticos por caminata, Tabla~\ref{tab:repetibilidad-fisica}. & Parcial; margen de estabilidad físico con el robot apoyado e IMU. \\
\bottomrule
\end{longtable}

La matriz viva, con criterio de cierre y rutas completas, se conserva en
\path{Documento_TESIS/MATRIZ_OBJETIVO_EVIDENCIA.md}.

\section{Protocolos aplicables}

Los procedimientos no se transcriben de forma abreviada para evitar dos
versiones divergentes. Las fuentes versionadas son:

\begin{itemize}
\item \path{Documentacion/PROTOCOLO_EXPERIMENTAL_BORRADOR.md}: ensayo, ciclo,
fallo, intervención, métricas y criterios provisionales;
\item \path{Documentacion/PROTOCOLO_CARACTERIZACION_FISICA.md}: fotografías,
geometría, masas y mapa de canales;
\item \path{Documentacion/PROTOCOLOS_FISICOS/}: puesta en marcha, calibración,
identificación y progresión de la marcha;
\item \path{Experimentos/PROTOCOLO_NODO_PRUEBAS_SEGURIDAD.md}: estímulos
aislados del supervisor;
\item \path{Experimentos/PROTOCOLO_PERTURBACIONES_GAZEBO.md}: condiciones para
fricción, empujes, ruido y retardos;
\item \path{Experimentos/PROTOCOLO_PPO_RESIDUAL.md}: contrato y campaña de la
capa aprendida;
\item \path{Experimentos/camaras/protocolo_repetibilidad.py}: ensayos físicos
con orden alternado, comprobación de salud y análisis tras cada caminata;
\item \path{Experimentos/camaras/ciclo_mejora.py}: ciclo de prueba en el robot,
actualización del modelo del actuador y reentrenamiento;
\item \path{Experimentos/camaras/calibrar_postura_colores.py}: ajuste de la
postura base del robot a la de la simulación con las cámaras.
\end{itemize}

\section{Lista de materiales}
\label{sec:lista-materiales}

La Tabla~\ref{tab:lista-materiales} recoge los componentes de la plataforma.
«En uso» indica que se empleó en los ensayos de octubre de 2026; «inventario»,
que está disponible en el laboratorio pero aún no se ha integrado; y «adquirir»,
que se propone comprarlo para completar la arquitectura de la
Figura~\ref{fig:esquema-propuesto}. Las especificaciones de los servomotores son
las de la ficha suministrada; en los ensayos se comportaron como servos de
posición de $180^\circ$.

\begin{longtable}{p{0.27\textwidth}p{0.05\textwidth}p{0.42\textwidth}p{0.13\textwidth}}
\caption{Lista de materiales de la plataforma y estado de cada componente.}\label{tab:lista-materiales}\\
\toprule
Componente & Cant. & Función y especificación & Estado \\
\midrule
\endfirsthead
\toprule
Componente & Cant. & Función y especificación & Estado \\
\midrule
\endhead
\multicolumn{4}{l}{\textit{Control}} \\
Raspberry~Pi~4 Model~B & 1 & Controlador embebido; Ubuntu~22.04, ROS~2, I$^2$C y GPIO a 3,3~V & en uso \\
Microcontrolador auxiliar & 1 & Interfaz de E/S por USB durante los ensayos de octubre & en uso \\
Controlador PWM PCA9685 & 1 & 16 canales, 12 bits, 50~Hz; dirección 0x40 & en uso \\
\multicolumn{4}{l}{\textit{Actuación}} \\
Servomotor MG996R & 12 & 4,8--6,6~V; 11~kgf$\cdot$cm a 6~V; 0,15~s/$60^\circ$; 55~g & en uso \\
\multicolumn{4}{l}{\textit{Sensores}} \\
Potenciómetro de eje & 10 & Realimentación de coxas, fémures y tibias 3 y 4 & en uso \\
Encoder magnético AS5600 & 12 & 12 bits, dirección 0x36; varios con imán débil & en uso \\
Multiplexor I$^2$C TCA9548A & 2 & Direcciones 0x70 y 0x71 para los AS5600 & en uso \\
Multiplexor analógico CD74HC4067 & 1 & 16 entradas para los potenciómetros & en uso \\
Convertidor ADS1115 & 1 & 16 bits, I$^2$C 0x48; lectura de los potenciómetros desde la Raspberry & adquirir \\
IMU BNO055 & 1 & Orientación del cuerpo (I$^2$C 0x28) para los ensayos apoyados & inventario \\
Monitor INA228 con derivador & 1 & Tensión, corriente y energía de los servos; dirección configurada distinta de 0x40 & adquirir \\
\multicolumn{4}{l}{\textit{Alimentación y seguridad}} \\
Fuente de los servos & 1 & 6~V regulada; corriente a dimensionar con el INA228 (orientativo: 20~A o más) & en uso \\
Fuente de la Raspberry & 1 & 5,1~V y 3~A, USB-C; se registró bajo voltaje con la actual & adquirir \\
Regulador LM2596 & 1 & Reductor para la electrónica auxiliar a 5~V & inventario \\
Fusible y portafusible & 1 & En la línea de los servos & adquirir \\
Parada de emergencia & 1 & Seta con contacto que corta la línea de los servos & adquirir \\
Bornera de masa común & 1 & Punto único de masa para servos y electrónica & adquirir \\
Conectores con traba & --- & JST-XH o similares para sensores y servos (los Dupont se soltaron en los ensayos) & adquirir \\
\multicolumn{4}{l}{\textit{Instrumentación de ensayo}} \\
Cámaras & 5 & Integrada del PC, Logitech C170, webcam USB y cámaras de dos portátiles & en uso \\
Marcadores ArUco y arnés & --- & Medición de la pose del cuerpo y sujeción en los ensayos apoyados & adquirir \\
\bottomrule
\end{longtable}

\section{Esquema eléctrico propuesto}

La Figura~\ref{fig:esquema-propuesto} propone el esquema de la plataforma con la
Raspberry~Pi como único controlador. Recoge las lecciones de los ensayos: la
electrónica y los servos tienen fuentes separadas y una masa común en un solo
punto; la línea de los servos lleva fusible, parada de emergencia y medición de
corriente; los potenciómetros se alimentan a 3,3~V para quedar dentro del rango
del ADS1115; y todos los dispositivos I$^2$C tienen direcciones distintas (el
INA228 debe configurarse fuera de 0x40, que es la del PCA9685). Los GPIO
indicados son una propuesta que debe confirmarse al cablear.

\begin{figure}[H]
\centering
\resizebox{\textwidth}{!}{%
\begin{tikzpicture}[
node distance=0.8cm and 1.1cm,
bl/.style={rectangle,draw,rounded corners,fill=blue!10,text width=3.1cm,minimum height=1.0cm,align=center,font=\footnotesize},
se/.style={rectangle,draw,rounded corners,fill=green!12,text width=3.1cm,minimum height=1.0cm,align=center,font=\footnotesize},
pw/.style={rectangle,draw,rounded corners,fill=orange!15,text width=3.1cm,minimum height=1.0cm,align=center,font=\footnotesize},
nu/.style={rectangle,draw,dashed,rounded corners,fill=gray!8,text width=3.1cm,minimum height=1.0cm,align=center,font=\footnotesize},
bus/.style={-,thick,blue!60!black},
pwr/.style={-latex,thick,orange!80!black},
sig/.style={-latex,semithick}]
\node[bl] (rpi) {Raspberry~Pi~4\\I$^2$C1: GPIO2 (SDA), GPIO3 (SCL)};
\node[bl,right=2.2cm of rpi] (pca) {PCA9685 (0x40)\\OE $\leftarrow$ GPIO4};
\node[bl,right=of pca] (servo) {12 servos MG996R};
\node[se,below=of pca] (tca) {TCA9548A 0x70 y 0x71};
\node[se,right=of tca] (as) {12 encoders AS5600 (0x36)};
\node[nu,below=of tca] (ads) {ADS1115 (0x48) \textit{(adquirir)}};
\node[se,right=of ads] (mux) {CD74HC4067\\S0--S3 $\leftarrow$ GPIO17, 27, 22, 23};
\node[se,right=of mux] (pot) {10 potenciómetros (3,3~V)};
\node[se,below=of ads] (imu) {IMU BNO055 (0x28)};
\node[nu,right=of imu] (ina) {INA228 \textit{(adquirir)}, dirección $\neq$ 0x40};
\node[pw,above=of servo] (fs) {Fuente servos 6~V};
\node[pw,above=of pca] (fus) {Fusible + parada de emergencia};
\node[pw,above=of rpi] (frpi) {Fuente 5,1~V / 3~A};
\node[pw,below=of rpi] (gnd) {Masa común (punto único: electrónica y servos)};
\draw[bus] (rpi.east) -- node[above,font=\scriptsize]{bus I$^2$C 3,3~V} (pca.west);
\draw[bus] ([xshift=1.1cm]rpi.east) |- (tca.west);
\draw[bus] ([xshift=1.1cm]rpi.east) |- (ads.west);
\draw[bus] ([xshift=1.1cm]rpi.east) |- (imu.west);
\draw[bus] (imu.east) -- (ina.west);
\draw[sig] (pca) -- node[above,font=\scriptsize]{PWM} (servo);
\draw[sig] (as) -- (tca);
\draw[sig] (pot) -- (mux);
\draw[sig] (mux) -- node[below,font=\scriptsize]{SIG$\rightarrow$A0} (ads);
\draw[pwr] (fs) -- (fus);
\draw[pwr] (fus.south) -- node[right,font=\scriptsize]{V+ (a través del INA228)} (fus.south |- pca.north);
\draw[pwr] (frpi) -- (rpi);
\draw[-,dashed] (gnd) -- (rpi);
\end{tikzpicture}}
\caption{Esquema eléctrico propuesto con la Raspberry~Pi como controlador
(bloques discontinuos: componentes por adquirir). Azul: bus I$^2$C; naranja:
alimentación; negro: señales.}
\label{fig:esquema-propuesto}
\end{figure}

\section{Comandos mínimos de reproducción}

La construcción y las pruebas deterministas se ejecutan desde la raíz:

{\footnotesize
\begin{verbatim}
source /opt/ros/humble/setup.bash
colcon build --packages-select \
  nova_sm3_description nova_gait_controller
source install/setup.bash
python3 -m pytest -q src/nova_gait_controller/test
\end{verbatim}
}

Los nueve escenarios dinámicos del supervisor se reproducen con:

{\footnotesize
\begin{verbatim}
ROS_DOMAIN_ID=93 \
Experimentos/ejecutar_pruebas_dinamicas_supervisor.sh \
  /tmp/pruebas_dinamicas_supervisor
\end{verbatim}
}

La línea base se lanza con \texttt{demo.launch.py}; antes de grabar se debe
comprobar instancia única mediante \texttt{ros2 node list | sort | uniq -c} y
confirmar los suscriptores del grabador en el tópico de órdenes. El análisis se
realiza con \path{Experimentos/analizar_gateo_rosbag.py} y
\path{Experimentos/analizar_contactos_rosbag.py}. La equivalencia se agrega con
\path{Experimentos/comparar_simuladores_equivalentes.py}; en MuJoCo se reinicia
primero el keyframe \texttt{stand}. Las rutas exactas de bolsa,
configuración y hashes se toman del registro de cada ensayo, nunca de una bolsa
con marca \texttt{ENSAYO\_INVALIDO.md}.

\section{Comandos de la fase física}

Entrenamiento de una semilla con el actuador real (ejemplo de gateo, desde
\path{Experimentos/entrenamiento_real_20261016}):

{\footnotesize
\begin{verbatim}
python -I entrenar_ppo_real.py --completo --aleatorizar \
  --marcha gateo --femur 15 --tibia 10 --apertura 8 \
  --timesteps 3000000 --lr 1e-4 --target-kl 0.03 --n-steps 2048 \
  --actuador actuadores_reales_v5.json --dt 0.03 \
  --w-avance 10 --w-actitud 30 --w-lateral 20 --w-giro 10 \
  --w-giro-ciclo 5 --w-suavidad 0.3 --w-coxas 1.0 \
  --seed 13 --output ppo_v14_gateo
\end{verbatim}
}

Selección de la mejor copia de una campaña y prueba en el robot (desde
\path{Experimentos/camaras}):

{\footnotesize
\begin{verbatim}
python -I ../entrenamiento_real_20261016/seleccionar_copias.py \
  --marcha gateo --carpeta ppo_v14_gateo --semillas 2
python3 protocolo_repetibilidad.py --sesion <nombre> --repeticiones 3 \
  --duracion-s 180 --pausa-s 120 \
  --condicion "gateo_base:gateo::--apertura 8" \
  --condicion "gateo_ppo:gateo:<política>/policy.zip"
python3 analizar_repetibilidad.py --sesion <nombre>
python3 ciclo_mejora.py --vueltas 2
\end{verbatim}
}

Las figuras de la fase física se regeneran con
\path{Documento_TESIS/generar_figuras_ppo_fisico.py}. Los videos no se versionan
por su tamaño; los registros de sensores, diagnósticos, análisis de video y
catálogos (\path{grabaciones/catalogo.csv} y \path{grabaciones/ciclos.csv}) sí.

\section{Reglas de integridad y repetición}

Cada resultado debe conservar fecha, versión, parámetros, marcadores de inicio
y fin, número de ciclos, criterio de exclusión y SHA-256. El ciclo 1 se mantiene
como transitorio; la unidad independiente es el ensayo. No se mezclan bolsas
creadas con cadencias, geometrías o perfiles diferentes. Las tentativas fallidas
se conservan con su causa y no entran en las estadísticas aceptadas.

El flujo de integración continua en
\path{.github/workflows/pruebas-automaticas.yml} compila los dos paquetes,
ejecuta las pruebas Python y repite los escenarios dinámicos en cada
\texttt{push} y \texttt{pull\_request}. Esto asegura regresión de software, no
reemplaza campañas dinámicas extensas ni validación física.

\section{Lista de transición al prototipo}

El orden F0--F4 se conserva en el plan de validación física fechado el 2 de
septiembre y versionado en la carpeta \path{Raspberry}. La
secuencia prevista es: ensamble y caracterización sin energía; seguridad
eléctrica y OE; calibración canal por canal; integración suspendida y luego con
soporte; instrumentación y contraste del modelo. En la fase de octubre se
ejecutaron la calibración, la integración suspendida y la instrumentación, y las
políticas PPO se ejecutaron solo con el robot suspendido, aunque la seguridad
eléctrica seguía abierta; esa condición la mitigaron los topes por
articulación, el límite de cambio por paso y la orden de apagado del firmware.
La integración con soporte y la marcha sobre el suelo requieren cerrar antes la
seguridad eléctrica y la alimentación.
\endgroup
